\documentclass[11pt,a4paper]{article}
\usepackage{fontspec,amsmath,amssymb,amsthm,microtype}
\setmainfont{lmroman10-regular.otf}[BoldFont=lmroman10-bold.otf,ItalicFont=lmroman10-italic.otf,BoldItalicFont=lmroman10-bolditalic.otf]
\setmonofont{lmmono10-regular.otf}[Scale=MatchLowercase]
\usepackage[margin=25mm]{geometry}
\usepackage[unicode,hidelinks]{hyperref}
\setlength{\emergencystretch}{3em}
\setlength{\parskip}{0.4em}
\setlength{\parindent}{0pt}
\allowdisplaybreaks[2]
\raggedbottom
\title{\texorpdfstring{The L\ensuremath{{}^{2}} splitting is the actual vector algorithm almost everywhere}{The L² splitting is the actual vector algorithm almost everywhere}}
\author{LangevinLean mathematical exposition}
\date{}
\begin{document}\maketitle
\noindent Source module: \texttt{Kinetic\allowbreak{}Langevin/\allowbreak{}Kinetic\allowbreak{}Splitting\allowbreak{}Lp.\allowbreak{}lean}. The accompanying HTML pairs this mathematical exposition with the exact Lean statements and records the scope of the formalization.


The deterministic convergence theorem applies to a Banach space of square-integrable random variables. To connect that analysis to the intended stochastic method, this module checks every elementary stage on actual representatives. The lifted force is the original force almost everywhere, and the supplied noise is retained unchanged. Induction then identifies the entire L\ensuremath{{}^{2}} recursion with the vector recursion.


\section{\texorpdfstring{Expose the chain driven by explicit filtered innovations}{Expose the chain driven by explicit filtered innovations}}


The innovation chain starts at its supplied phase state and recursively applies the explicit Verlet–OU step with innovation \(N_k\). The earlier continuous-input chain is exactly this chain with \(N_k=N_{kh,h}(B)\). Lean proves this equality by induction, unfolding the two actual recursions at the successor step. This representation lets the probability layer identify the innovation array law independently of the deterministic input analysis.



{\raggedright\noindent Lean declarations: \texttt{kinetic\allowbreak{}Split\allowbreak{}Noise\allowbreak{}Chain}, \texttt{kinetic\allowbreak{}Split\allowbreak{}Chain\_\allowbreak{}eq\_\allowbreak{}noise\allowbreak{}Chain}.\par}

\section{\texorpdfstring{Check force kicks on actual L\ensuremath{{}^{2}} representatives}{Check force kicks on actual L² representatives}}


The lifted force \(G\) is constructed from the actual Lipschitz force \(g\), and the shared lifting theorem gives \(G(X)(\omega)=g(X(\omega))\) almost everywhere. The representative identities for subtraction and scalar multiplication therefore give \([V-hG(X)](\omega)=V(\omega)-hg(X(\omega))\). Position is unchanged. Lean intersects the three almost-everywhere events and proves equality of the two phase coordinates. No new vector-field or output-law assumption is introduced.



{\raggedright\noindent Lean declarations: \texttt{lp\_\allowbreak{}kick\allowbreak{}Step\_\allowbreak{}coe}.\par}

\section{\texorpdfstring{Check free transport and the full supplied thermostat noise}{Check free transport and the full supplied thermostat noise}}


The representative of \(X+hV\) is \(X(\omega)+hV(\omega)\), proving the free drift identity. Likewise the representative of \(\delta V+\sigma\Xi\) is \(\delta V(\omega)+\sigma\Xi(\omega)\), proving the OU stage identity. Lean explicitly intersects the addition and scalar-multiplication representative events and retains the exact supplied innovation \(\Xi\). These are pathwise stage identities and do not presume any Gaussian distribution.



{\raggedright\noindent Lean declarations: \texttt{lp\_\allowbreak{}drift\allowbreak{}Step\_\allowbreak{}coe}, \texttt{lp\_\allowbreak{}ornstein\allowbreak{}Step\_\allowbreak{}coe}.\par}

\section{\texorpdfstring{Compose the checked stages into the actual Verlet–OU map}{Compose the checked stages into the actual Verlet–OU map}}


The Verlet step is the kick–drift–kick composition. Apply the checked representative identity at each actual intermediate L\ensuremath{{}^{2}} state, intersect the resulting almost-everywhere events, and substitute them in order. The complete splitting then composes this checked Verlet map with the checked OU stage, using exactly the damping \(e^{-\gamma h}\) and supplied filtered innovation. Lean performs these successive substitutions, so the state at which the second force is evaluated remains the updated position.



{\raggedright\noindent Lean declarations: \texttt{lp\_\allowbreak{}verlet\allowbreak{}Step\_\allowbreak{}coe}, \texttt{lp\_\allowbreak{}kinetic\allowbreak{}Split\allowbreak{}Step\_\allowbreak{}coe}.\par}

\section{\texorpdfstring{Identify every numerical iterate with the original vector chain}{Identify every numerical iterate with the original vector chain}}


The initial phase representatives agree by definition. Assuming agreement at step \(n\), the checked full-step identity and the induction hypothesis give agreement at step \(n+1\), driven by the same innovation representative \(N_n(\omega)\). Lean intersects these two almost-everywhere events and rewrites the recursive definitions. Hence every L\ensuremath{{}^{2}} iterate analyzed by the convergence theorem is the actual vector splitting algorithm on the same probability source.



{\raggedright\noindent Lean declarations: \texttt{lp\_\allowbreak{}kinetic\allowbreak{}Split\allowbreak{}Noise\allowbreak{}Chain\_\allowbreak{}coe}.\par}

\end{document}
